\documentclass[10pt,a4paper]{amsart}
\usepackage[margin=19mm]{geometry}
\usepackage{amsmath,amssymb,mathtools}
\usepackage[scr=rsfs,cal=euler]{mathalfa}
\usepackage{xfrac}
\usepackage[unicode,hidelinks,bookmarks=false]{hyperref}
\newcommand{\ie}{i.e.\ }
\newcommand{\cf}{cf.\ }
\newcommand{\resp}{respectively\ }
\newcommand{\Subsection}{\subsection}
\providecommand{\nobreakdash}{\nobreak\hbox{-}}
\setlength{\emergencystretch}{2em}
\multlinegap0pt
\usepackage{tikz}
\usetikzlibrary{cd}
\usepackage{mathtools}
\usepackage{amssymb}
\usepackage{tikz}
\newtheorem{prop}{Proposition}
\newcommand{\on}{\operatorname}
\title{Comparison of Motivic Homotopy Theories}
\author{Tianjian Tan}
\date{Source: arXiv:2603.11753v1 (CC BY 4.0)}
\begin{document}
\maketitle

\noindent\emph{Source and licence.} Comparison of Motivic Homotopy Theories. Version arXiv:2603.11753v1 (CC BY 4.0), distributed by arXiv under the Creative Commons Attribution 4.0 International licence, http://creativecommons.org/licenses/by/4.0/, as recorded in the arXiv licence metadata for this exact version and shown on the abstract page https://arxiv.org/abs/2603.11753v1. Full licence notice: \texttt{licenses/arxiv-2603.11753-CC-BY-4.0.txt}.\par\smallskip

\noindent\textit{Editorial scope.} Counted unit: the article's prop ver. (source0.tex lines 1053--1061). The article's complete original proof of it is delivered with it (source0.tex lines 1062--1132, one \texttt{proof} environment of 71 lines). No mathematics is written, repaired or re-derived: the statement and the proof are the article's own lines, in the article's own notation, and no retained line is substituted or re-flowed.  Retained uncounted as the source-local context the proof needs: lines 1039--1052.  Nothing was omitted from any retained span, so the package carries no omission record.  No retained line contains a citation, so the wrapper carries no bibliography and no bibliography entry.  The retained lines contain the article's own diagram code, which the wrapper typesets with the standard tikz packages; no image file is delivered and the package declares no asset dependency.  Editorial formatting only, in addition: an AMS \texttt{amsart} preamble that restates the article's own declarations and macros used by the retained lines, namely the environments \texttt{prop} on separate counters, together with the standard packages those lines need.  The retained environments print in this document as Proposition 1 in order of appearance, whereas the article prints the same items with the numbering of its own full document.  The complete native source of the article as distributed in the arXiv e-print for this exact version is bundled unchanged as \texttt{sources/196235/source0.tex}.
\begin{prop}\label{cons}
The composition $\on{L}_{(\mathbb{P}^1)^{dual}}\circ\on{L}_{\on{coNis},\on{co}{\text -}\mathbb{A}^1}\circ j$ induces, for each presentable, stable category $\mathcal{D}$, a fully faithful functor\[
\on{Fun}^L(\mathcal{SH}(S)^\vee,\mathcal{D})\rightarrow\on{Fun}(Sm^{\on{fp}}(S)^{op},\mathcal{D})
\]
with the essential image spanned by $\mathbb{A}^1$-invariant functors $F$ that preserve terminal object, send Nisnevich squares to pushouts, and send the fiber\[fib(F(\mathbb{P}^1)\xrightarrow{F(\infty)}F(S))\]
to an invertible object.


Moreover, since the symmetric monoidal structures are compatible, the claim remains valid for symmetric monoidal category $\mathcal{D}$ and symmetric monoidal functors.


In particular, we have\[\on{Fun}^{L,\otimes}(\mathcal{SH}(S)^\vee,\on{Mot}(S)_{\on{loc}}^{\mathbb{A}^1})\simeq\on{Fun}^{\otimes}_{\on{Nis},\mathbb{A}^1,\mathbb{P}^1}(Sm^{\on{fp}}(S)^{op},\on{Mot}(S)_{\on{loc}}^{\mathbb{A}^1})
\]
\end{prop}

\begin{prop}\label{ver.}
The functor\[
\begin{aligned}
\phi: Sm^{\on{fp}}(S)^{op} &\rightarrow \on{Mot}(S)_{\on{loc}}^{\mathbb{A}^1}\\
 X & \mapsto \mathcal{U}(\on{perf}_X)
 \end{aligned}\]
satisfies all the conditions in the Proposition \ref{cons}. In particular, it extends to a symmetric monoidal left adjoint functor \[ \Phi:\mathcal{SH}(S)^\vee \rightarrow \on{Mot}(S)_{\on{loc}}^{\mathbb{A}^1}.
\]
\end{prop}

\begin{proof}
\begin{itemize}
\item Compatibility of symmetric monoidal structures.\\
Both of the $S$-linear functors\[
\begin{gathered}
\on{perf}:Sm(S)^{op}\rightarrow \mathsf{Cat}^{\on{perf}}(\on{perf}_S)\\
\mathcal{U}:\mathsf{Cat}^{\on{perf}}(\on{perf}_S)\rightarrow\on{Mot}(S)_{\on{loc}}^{\mathbb{A}^1}
\end{gathered}
\]
are symmetric monoidal. So is their composition.
\item Nisnevich excision.\\
It is obvious that the empty scheme is sent to $\mathcal{U}(\on{perf}_{\emptyset})=0$. On the other hand, recall that the functor\[
\on{perf}:Sm(S)^{op}\rightarrow \mathsf{Cat}^{\on{perf}}(\on{perf}_S)\]
satisfies fppf descent. In particular, for any Nisvevich square
\[\begin{tikzcd}
	W && V \\
	\\
	U && X
	\arrow[hook, from=1-1, to=1-3]
	\arrow[from=1-1, to=3-1]
	\arrow[from=1-3, to=3-3]
	\arrow[hook, from=3-1, to=3-3],
\end{tikzcd}\]
we have a pullback diagram
\[\begin{tikzcd}
	{\on{perf}_X} && {\on{perf}_U} \\
	\\
	{\on{perf}_V} && {\on{perf}_W}
	\arrow[from=1-1, to=1-3]
	\arrow[from=1-1, to=3-1]
	\arrow[from=1-3, to=3-3]
	\arrow[from=3-1, to=3-3].
\end{tikzcd}\]
For an open immersion $U\hookrightarrow X$, if we denote by $\on{perf}_{X,X-U}$ the fiber\[
\on{perf}_{X,X-U}\rightarrow\on{perf}_X\rightarrow\on{perf}_U,\]
then we have a commutative diagram
\[\begin{tikzcd}
	{\on{perf}_{X,X-U}} && {\on{perf}_X} && {\on{perf}_U} \\
	\\
	{\on{perf}_{V,V-W}} && {\on{perf}_V} && {\on{perf}_W}
	\arrow[from=1-1, to=1-3]
	\arrow["\simeq"', from=1-1, to=3-1]
	\arrow[from=1-3, to=1-5]
	\arrow[from=1-3, to=3-3]
	\arrow[from=1-5, to=3-5]
	\arrow[from=3-1, to=3-3]
	\arrow[from=3-3, to=3-5].
\end{tikzcd}\]
Since $\mathcal{U}$ is the universal localizing invariant, applying $\mathcal{U}$ gives a commutative diagram in which, again, the horizontal sequences are fibers and the left most arrow is an equivalence. As a consequence,
\[\begin{tikzcd}
	{\mathcal{U}(\on{perf}_X)} && {\mathcal{U}(\on{perf}_U)} \\
	\\
	{\mathcal{U}(\on{perf}_V)} && {\mathcal{U}(\on{perf}_W)}
	\arrow[from=1-1, to=1-3]
	\arrow[from=1-1, to=3-1]
	\arrow[from=1-3, to=3-3]
	\arrow[from=3-1, to=3-3]
\end{tikzcd}\]
is a pullback as desired.
\item $\mathbb{A}^1$-invariance.\\
This is true by impose.
\item $\mathbb{P}^1$-invertibility.\\
We have a cofiber sequence:
\[\begin{tikzcd}
	{\on{perf}_S} && {\on{perf}_{\mathbb{P}^1}} && {\on{perf}_S}
	\arrow[from=1-1, to=1-3]
	\arrow["\on{perf}_\infty", from=1-3, to=1-5].
\end{tikzcd}\]
And $\mathcal{U}$ sends it to a both fiber and cofiber sequence. In particular, the fiber of \[\phi(\mathbb{P}^1\xrightarrow{\infty} S)\] is the unit, thus is invertible.
\end{itemize}
\end{proof}

\end{document}
